\section{Conclusion}
\label{sec:conclusion}
This paper develops a tractable principal--agent framework to study agency conflicts in delegated portfolio management. In the model, performance chasing generates an option-like payoff for managers, which tilts their incentives toward higher volatility relative to what investors would choose. The central contribution of the paper is to characterize how compensation contracts can be designed to implement a desired level of portfolio risk.

The main practical output is a fee schedule that implements a target volatility: for each risk target, there exists a management fee level such that the manager's first-order condition is satisfied at that target, making it a locally incentive-compatible equilibrium under the stated conditions (global uniqueness is verified numerically across the calibrated parameter grid). Rather than increasing volatility to benefit from convex flow dynamics, the manager chooses the investor's preferred target. We also show that a simple linear approximation to the exact incentive-compatible schedule is accurate over empirically relevant ranges, making the contract easy to communicate and implement as ``pay for risk, not luck.'' Finally, the numerical exercise illustrates how flow convexity, fund scale, managerial skill, and risk aversion shape the equilibrium risk target and the fee level required to align incentives.

A natural next step in this research is to incorporate run-like dynamics and risk of ruin: when redemptions become strategic and non-linear in bad states, the same incentive problem remains, but the manager's costs of excessive risk can become impossible to disregard as AUM falls in a concave fashion when losses are large.
